\section{corrispondenza tra integrali di superficie e integrali di forme differenziali}

In questo capitolo mettiamo in corrispondenza gli integrali di superificie con gli integrali 
di forme differenziali nei casi più rilevanti, ovvero per $k=1$ (curve), per $k=n-1$ (ipersuperfici) 
e per $k=n$ (volumi). 
Nel caso rilevante $n=3$ si esaurisce tutta la casistica, lo tratteremo in particolare.

\subsection{curve}
Per le curve abbiamo già visto che se $\xi$ è un campo vettoriale in $\RR^n$ 
possiamo associare a $\xi$ la $1$-forma differenziale $\omega$ definita da
\[
  \omega_x[v] = \xi(x) \cdot v.
\]
In tal modo si ha
\[
   \int_\phi \omega = \int_\phi \xi(x)\cdot \phi'(t)\, dt.
\]
Se $\phi'(t) \neq 0$ si può definire $\tau_\phi(t) \defeq \phi'(t)/\abs{\phi'(t)}$ 
che geometricamente è il versore tangente alla curva rappresentata da $\phi$ e dunque
si ha 
\[
  \phi'(t)\, dt = \tau_\phi \cdot \abs{\phi'(t)}\, dt = \tau_\phi\, ds
\]
dove $ds$ è l'elemento di lunghezza della curva (ovvero $s$ è il parametro di lunghezza d'arco della curva).
Osserviamo che sebbene $\tau_\phi$ non ha senso quando $\phi'(t) = 0$ il prodotto
$\tau_\phi \abs{\phi'(t)} = \phi'(t)$ è sempre ben definito, dunque gli integrali 
della forma 
\[
  \int_\phi \xi \cdot \tau_\phi\, ds
\] 
hanno senso anche se $\phi'(t)$ si annulla.

\subsection{ipersuperfici}

Ad un campo vettoriale $\xi\colon A\subset \RR^n\to \RR^n$ possiamo associare la 
$(n-1)$-forma differenziale $\omega$ definita da
\[
  \omega_x[v_2,\dots,v_n] = \det(\xi(x), v_2,\dots,v_n).
\]
In tal modo, grazie a~\eqref{eq:493444} si ottiene che 
l'integrale di $\omega$ su una ipersuperficie rappresentata da $\phi\colon B\subset \RR^{n-1} \to \RR^n$
coincide con il flusso del campo vettoriale $\xi$ attraverso la ipersuperficie rappresentata da $\phi$: 
\begin{align*}
  \int_\phi \omega
  &= \int_B \det (\xi(x), D\phi(x))\, dx \\
  &= \int_B \xi(x) \cdot \nu_\phi(x) \, \sqrt{\det((D\phi)^t D\phi)}\, dx \\
  &= \int_\phi \xi \cdot \nu_\phi\, d\sigma.
\end{align*}

Se $S = \current{\phi_1} + \dots + \current{\phi_M}$ è una $(n-1)$-superficie parametrizzata 
ha quindi senso definire 
\[
   \int_S \xi \cdot \nu_S\, d\sigma \defeq \int_S \omega
\]
dove $\omega$ è la $(n-1)$-forma differenziale associata a $\xi$ come sopra.

\subsection{volumi}

Se $A\subset \RR^n$ è un aperto limitato e misurabile si può definire $\current A\colon \Omega^n(A) \to \RR$ tramite
\[
  \int_{\current{A}} f \, dx_1 \wedge \dots \wedge dx_n 
  = \int_A f.
\]
Se $A$ è sufficientemente regolare sarà possibile scrivere 
$\current A = \current \phi_1 + \dots + \current \phi_M$ 
con $\phi_i\colon [0,1]^n \to A$ di classe $C^1$. 
In tal caso $\current A$ è una $n$-superficie parametrizzata e dunque si potrà definire il suo 
bordo. 

\begin{theorem}
Sia $A\subset \RR^n$ un aperto limitato sufficientemente regolare (in modo che $\current A$ sia una $n$-superficie parametrizzata).
Si $f\colon \bar A \to \RR^n$ una funzione di classe $C^1$.
Allora 
\[
   \int_A \div f = \int_{\partial \current A} f \cdot \nu_{\partial \current A}\, d\sigma.
\]
\end{theorem}
\begin{proof}
Consideriamo la $(n-1)$-forma $\omega$ associata a $f$.
Il teorema~\ref{th:stokes} ci dice che 
\[
  \int_{A} \div f
  = \int_{\current A} d\omega 
  = \int_{\current A} \omega
  = \int_{\partial \current A} f \cdot \nu_{\partial \current A}\, d\sigma.
\]
\end{proof}

\subsection{forme differenziali multilineari in $\RR^3$}

\begin{definition}[prodotto vettore]
  \label{def:prodotto_vettore}
Siano $v,w\in \RR^3$.
Definiamo il \emph{prodotto vettore} $v\times w\in \RR^3$ 
come l'unico vettore tale che per ogni $u\in \RR^3$ si abbia
\[
  u \cdot (v\times w) = \det\enclose{u, v, w} \qquad\text{(prodotto triplo)}.
\]
Questa uguaglianza definisce univocamente $v\times w$ in quanto 
il lato destro è una applicazione lineare in $u$ e dunque 
c'è un unico vettore che la rappresenta mediante prodotto scalare.
\end{definition}

Ovviamente il vettore $v\times w$ è perpendicolare sia a $v$ che a $w$,
perché $v\cdot (v\times w) = \det\enclose{v,v,w} = 0$ e $w\cdot (v\times w) = \det\enclose{w,v,w} = 0$.
Posto $\nu = \frac{v\times w}{\abs{v\times w}}$ si ha che $\nu$ è un versore normale al piano
generato da $v$ e $w$.
Inoltre 
\[ 
\abs{v\times w} = \nu\cdot(v\times w) = \det\enclose{\nu, v, w}
\]
è il volume del prisma retto generato da $v$ e $w$ in direzione $\nu$
ed essendo $\nu$ unitario è anche l'area del parallelogramma generato da $v$ e $w$ in $\RR^3$.
Significa che $v\times w$ è un vettore con modulo pari all'area del parallelogramma generato da $v$ e $w$ 
e con direzione $\nu$ perpendicolare a $v$ e $w$ secondo la regola della mano destra
(in modo che $\det \enclose{v, w, \nu} > 0$).
Se $u$ e $v$ sono tra loro paralleli, si ha $u\times v = 0$.

In coordinate, se $\xi = v\times w$ con $v=(v_1,v_2,v_3)$, $w=(w_1,w_2,w_3)$
si ha dunque $\xi=(\xi_1,\xi_2,\xi_3)$ 
con
\begin{align*}
  \xi_1 &= \vec e_1 \cdot \xi = \det\enclose{\vec e_1,v,w} = v_2 w_3 - v_3 w_2,\\ 
  \xi_2 &= \vec e_2 \cdot \xi = \det\enclose{\vec e_2,v,w} = v_3 w_1 - v_1 w_3,\\
  \xi_3 &= \vec e_3 \cdot \xi = \det\enclose{\vec e_3,v,w} = v_1 w_2 - v_2 w_1.
\end{align*}

Più espressivamente si può scrivere
\[
  \xi = \det\begin{pmatrix}
    \vec e_1 & v_1 & w_1 \\ \vec e_2 & v_2 & w_2 \\ \vec e_3 & v_3 & w_3
  \end{pmatrix}
  = (v_2 w_3 - v_3 w_2)\, \vec e_1 + (v_3 w_1 - v_1 w_3)\,  \vec e_2 + (v_1 w_2 - v_2 w_1)\, \vec e_3. 
\]

Se $\omega\in \Omega^2(A)$ è rappresentata dal campo vettoriale 
$\xi\colon A\to \RR^3$ 
e $\phi\colon B\subset \RR^2\to A \subset \RR^3$ 
osserviamo che $\phi_x = \frac{\partial \phi}{\partial x}$ e $\phi_y = \frac{\partial \phi}{\partial y}$ 
sono due vettori tangenti alla superficie parametrizzata da $\phi$ e dunque 
\[
  \nu_\phi = \frac{\phi_x \times \phi_y}{\abs{\phi_x \times \phi_y}}
\]
è il versore normale alla superficie in $\phi(x,y)$ 
con direzione che dipende dall'orientazione di $\phi$.
L'integrale di $\omega$ lungo la superficie parametrizzata da $\phi$ si 
scrive dunque come
\begin{align*}
 \int_{\phi} \omega
 &= \int_B \omega \Enclose{\phi_x, \phi_y}\, dx\, dy
 = \int_B \det \enclose{\xi, \phi_x, \phi_y}\, dx\, dy\\
 &= \int_B \xi \cdot \enclose{\phi_x \times \phi_y}\, dx\, dy
 = \int_B \xi \cdot \nu_\phi\, d\sigma
\end{align*}
dove $d\sigma = \abs{\phi_x \times \phi_y}\, dx\, dy$ 
rappresenta l'elemento di area della superficie parametrizzata da $\phi$.
Dunque $\int_\phi \omega$ è il flusso del campo vettoriale $\xi$ 
attraverso la superficie parametrizzata da $\phi$.

\begin{definition}[rotore]
Sia $A$ un aperto di $\RR^3$ e $\xi \colon A \to \RR^3$ una funzione differenziabile.
Sia $\vec e_1,\vec e_2,\vec e_3$ la base canonica di $\RR^3$.
Definiamo allora il \emph{rotore} di $\xi$ come il vettore $\rot \xi \in \RR^3$ 
definito da 
\[
   \rot \xi = \nabla \times \xi
    = \det \begin{pmatrix}
    \vec e_1 & \partial_x & \xi_1 \\
    \vec e_2 & \partial_y & \xi_2 \\
    \vec e_3 & \partial_z & \xi_3 
    \end{pmatrix}
    = \begin{pmatrix}
    \frac{\partial \xi_3}{\partial y} - \frac{\partial \xi_2}{\partial z} \smallskip \\
    \frac{\partial \xi_1}{\partial_z} - \frac{\partial \xi_3}{\partial x} \smallskip \\
    \frac{\partial \xi_2}{\partial x} - \frac{\partial \xi_1}{\partial y}
    \end{pmatrix}.
\]

Ricordiamo anche gli operatori \emph{gradiente} e \emph{divergenza}
\[
  \grad f = \nabla f = \begin{pmatrix}
     \frac{\partial f}{\partial x} \smallskip \\
     \frac{\partial f}{\partial y} \smallskip \\
     \frac{\partial f}{\partial z}
  \end{pmatrix}, \qquad 
  \div f = \nabla \cdot \xi = \frac{\partial \xi_1}{\partial x} + \frac{\partial \xi_2}{\partial y} + \frac{\partial \xi_3}{\partial z}.
\]
In tutte queste notazioni abbiamo usato l'operatore \emph{nabla}: $\nabla = (\partial_x, \partial_y, \partial z)$
dove $\partial_x = \frac{\partial}{\partial x}$, $\partial y = \frac{\partial}{\partial y}$, $\partial_z = \frac{\partial}{\partial z}$:
\end{definition}

Abbiamo già visto che una $1$-forma $\omega$ in un aperto $A$ di $\RR^3$ può essere rappresentata 
con un campo vettoriale $\xi\colon A \subset \RR^3 \to \RR^3$ 
in modo tale che si abbia 
\[
  \omega_{(x,y,z)}[v] = \xi(x,y,z)\cdot v.
\]

Se consideriamo $f\colon A\subset \RR^3\to \RR$ una $0$-forma, 
il suo differenziale $\omega = df$ è una $1$-forma, che si rappresenta 
con il vettore $\nabla f$:
\[
df = \frac{\partial f}{\partial x}\, dx 
+ \frac{\partial f}{\partial y} dy\, \frac{\partial f}{\partial z}\, dz.
\]

Se $\omega\in \Omega^1(A)$, possiamo scrivere $\omega$ nella forma 
\[
 \omega = g_1\, dx + g_2\, dy + g_3\, dz
\]
in modo che $\vec g \colon A\to\RR^3$ sia il campo vettoriale che rappresenta $\omega$.
In questo caso si ha 
\begin{align*}
    d\omega &= (\partial_x \vec g_1 dx + \partial_y \vec g_1 dy + \partial_z \vec g_3 dz)\wedge dx 
    + (\partial_x \vec g_2 dx + \partial_y \vec g_2 dy + \partial_z \vec g_3 dz)\wedge dy \\
    & \quad + (\partial_x \vec g_3 dx + \partial_y \vec g_3 dy + \partial_z \vec g_3 dz)\wedge dz \\    
    &= (\partial_y \vec g_3 - \partial_z \vec g_2)\, dy \wedge dz
       + (\partial_z \vec g_1 - \partial_x \vec g_3)\, dz\wedge dx 
       + (\partial_x \vec g_2 - \partial_y \vec g_1)\, dx \wedge dy\\
    &= (\rot \vec g)_1 \, dy \wedge dz
       + (\rot \vec g)_2\, dz\wedge dx 
       + (\rot \vec g)_3\, dx \wedge dy.
\end{align*}
Dunque se rappresentiamo $d\omega\in \Omega^2(A)$, tramite 
il campo vettoriale $\xi\colon A\to \RR^3$
definito da 
\[
\xi_1(\vec p) = \omega_{\vec p}[\vec e_2,\vec e_3], \qquad
\xi_2(\vec p) = \omega_{\vec p}[\vec e_3,\vec e_1], \qquad
\xi_3(\vec p) = \omega_{\vec p}[\vec e_1,\vec e_2]
\]
si avrà
\begin{equation}\label{eq:fkjtpd}
  d \omega = \xi_1\, dy \wedge dz + \xi_2\, dz\wedge dx + \xi_3\, dx\wedge dy.
\end{equation}
Abbiamo quindi verificato che se $\xi$ è il campo vettoriale che rappresenta 
la $1$-forma $\omega$, il campo (pseudo)-vettoriale che rappresenta $d\omega$
è $\rot \xi$.

Consideriamo infine $\omega \in \Omega^2(A)$ una $2$-forma 
definita su un aperto $A\subset\RR^3$.
Sappiamo che $\omega$ si rappresenta tramite un campo
vettoriale $\xi \colon A \subset \RR^3 \to \RR^3$  
come in~\eqref{eq:fkjtpd}.
Osserviamo che allora si ha 
\begin{align*}
  \omega[v,w] 
  &= \xi_1 (v_1 w_3 - v_3 w_2) + \xi_2 (v_3 w_1 - v_1 w_3)
   + \xi_3 (v_1 w_2 - v_2 w_1) \\
   &= \det \enclose{\xi, v, w} = \xi \cdot (v\times w).
\end{align*}
Inoltre
\begin{align*}
 d\omega &= (\partial_x \xi_1 dx + \partial_y \xi_1 dy + \partial_z \xi_1 dz) \wedge dy \wedge dz \\
    &\quad + (\partial_x \xi_2 dx + \partial_y \xi_2 dy + \partial_z \xi_2 dz) \wedge dz \wedge dx \\
    &\quad + (\partial_x \xi_3 dx + \partial_y \xi_3 dy + \partial_z \xi_3 dz) \wedge dx \wedge dy \\
    &= (\partial_x \xi_1 + \partial_y \xi_2 + \partial_z \xi_3)\, dx \wedge dy \wedge dz\\
    &= (\div \xi)\, dx \wedge dy \wedge dz.
\end{align*}
Dunque la funzione che rappresenta la $3$-forma $d\omega$ 
è la divergenza del campo vettoriale che rappresenta la $2$-forma $\omega$.

Possiamo quindi esprimere il teorema~\ref{th:stokes} in termini di rotore.
\begin{theorem}[del rotore]
Sia $\xi\colon A\to \RR^3$ un campo vettoriale di classe $C^1$.
Allora si ha 
\[
 \int_S \rot \xi = \int_{\partial S} \xi
\]

\end{theorem}